\section{Methodology}

\subsection{Proposed Approach}
This study proposes a computer vision-based approach for detecting traffic signal violations using CCTV footage from a static camera installed at Simpang Wirosaban, Yogyakarta. The system is designed as a multi-stage pipeline (Figure \ref{fig:diagram_process}) consisting of image enhancement, vehicle detection, multi-object tracking, traffic light recognition, and violation analysis. To reduce computational overhead and eliminate detections from irrelevant regions such as sidewalks and surrounding infrastructure, a polygon-based Region of Interest (ROI) is defined to constrain the analysis strictly to the active roadway area.

Prior to vehicle detection, each frame undergoes a three-stage image enhancement pipeline to improve visibility under nighttime and uneven lighting conditions. Gamma correction is first applied to brighten underexposed regions by adjusting pixel intensity according to a tunable gamma value ($\gamma = 1.5$). The gamma-corrected frame is then passed to Contrast Limited Adaptive Histogram Equalization (CLAHE), applied exclusively to the luminance channel of the LAB color space to enhance local contrast without introducing chrominance distortion. Finally, Unsharp Masking is applied to sharpen object boundaries by subtracting a Gaussian-blurred version of the frame from the original, which mitigates the effect of motion blur on moving vehicles. This enhancement pipeline is applied solely to the input of the vehicle detector; the original unmodified frame is retained for all other processing stages to preserve color fidelity.

Within the defined ROI, vehicle detection is performed using YOLO11, which identifies vehicles present inside the polygonal area on each frame. The detected bounding boxes are subsequently processed by the Simple Online and Realtime Tracking (SORT) algorithm, which employs a Kalman filter for state prediction and the Hungarian algorithm for frame-to-frame identity association. This tracking mechanism assigns a persistent unique identifier to each vehicle and maintains its trajectory across consecutive frames, enabling reliable monitoring of individual vehicle movements over time.

In parallel with vehicle detection, the traffic light state is determined using an HSV color-based thresholding method applied to a dedicated ROI positioned at the traffic light location in the frame. The detected state is used as a conditional input to the violation logic module. A virtual stop line is defined within the ROI to represent the legal stopping boundary at the intersection. Violation detection is triggered when a tracked vehicle, whose trajectory originated from within the ROI, crosses the stop line at the moment the traffic light is in a red state. To prevent duplicate records, each unique vehicle identity is registered as a violation at most once. The system produces output in the form of annotated video frames displaying bounding boxes, vehicle identities, stop line geometry, and a running violation counter, providing a structured and automated means of traffic violation documentation.

\subsection{Dataset}
Write the dataset description here.

\subsection{Image Enhancement}
Write the preprocessing steps here.

\subsection{Detection and Tracking}
Write the detection and tracking method here.

\subsection{Violation Logic}
Write the violation decision logic here.